% Part: history
% Chapter: biographies 
% Section: alfred-tarski

\documentclass[../../../include/open-logic-section]{subfiles}

\begin{document}

\olfileid{his}{bio}{tar}

\olsection{આલ્ફ્રેડ ટાર્સ્કી}

\olphoto{tarski-alfred}{આલ્ફ્રેડ ટાર્સ્કી}

આલ્ફ્રેડ ટાર્સ્કીનો જન્મ 14 જાન્યુઆરી
1901ના રોજ પોલૅન્ડના વૉર્સોમાં થયો હતો,
જે ત્યારે રશિયન સામ્રાજ્યનો ભાગ હતું.
“નેપોલિયન જેવા” વર્ણવાયેલા ટાર્સ્કી
ઉત્સાહી, વાતોડિયા અને તીવ્ર સ્વભાવના
હતા. તેમની શક્તિ ઘણી વાર તેમના
વ્યાખ્યાનોમાં પણ દેખાતી: એક વાર
વ્યાખ્યાન દરમિયાન સિગારેટ ફેંકતાં
તેમણે કચરાપેટીમાં આગ લગાડી દીધી
અને પછી તેમને એ ઇમારતમાં
ફરી વ્યાખ્યાન આપવાની મનાઈ થઈ.

ટાર્સ્કીને નાનપણથી જ જ્ઞાનની
તરસ હતી. જીવનમાં પાછળથી
તેઓ વિદ્યાર્થીઓને કહેતા કે
તર્કશાસ્ત્ર એકમાત્ર એવો વર્ગ
હતો જેમાં તેમને બી શ્રેણી
મળી હતી, એટલે તેમણે
તર્કશાસ્ત્ર ભણ્યું. પરંતુ
માધ્યમિક શાળાના તેમના
અભિલેખ બતાવે છે કે
તર્કશાસ્ત્ર સહિત દરેક
વિષયમાં તેમને એ શ્રેણી
મળી હતી. તેમણે 1918થી
1924 સુધી વૉર્સો
યુનિવર્સિટીમાં અભ્યાસ
કર્યો. શરૂઆતમાં જીવવિજ્ઞાન
ભણવાનો વિચાર હતો, પરંતુ
યુનિવર્સિટી વૉર્સોની
તર્કશાસ્ત્ર અને તત્વજ્ઞાનની
શાળાનું કેન્દ્ર હોવાથી
તેમને ગણિત, તત્વજ્ઞાન અને
તર્કશાસ્ત્રમાં રસ પડ્યો.
1924માં સ્તાનિસ્વાવ
લેશ્નેવ્સ્કીના માર્ગદર્શન
હેઠળ તેમણે ડૉક્ટરેટ મેળવી.

1939માં અમેરિકામાં સ્થળાંતર
કરતાં પહેલાં ટાર્સ્કીએ
વૉર્સોમાં માધ્યમિક શાળાના
શિક્ષક તરીકે કામ કરતા
પોતાનું કેટલાક સૌથી મહત્વનું
સંશોધન પૂરું કર્યું.
તાર્કિક ફલિતતા અને
તાર્કિક સત્ય વિશેના
તેમના લેખો એ જ સમયમાં
લખાયા. 1939માં તેઓ
વ્યાખ્યાન પ્રવાસે
અમેરિકા ગયા હતા. ત્યાં
તેમના રોકાણ દરમિયાન
જર્મનીએ પોલૅન્ડ પર
આક્રમણ કર્યું; યહૂદી
વંશના હોવાથી તેઓ
પાછા જઈ શક્યા નહીં.
તેમનાં પત્ની અને
બાળકો યુદ્ધ પૂરું થાય
ત્યાં સુધી પોલૅન્ડમાં
રહ્યાં, પરંતુ ત્યાર
પછી તેઓ પણ અમેરિકામાં
સ્થળાંતર કરી શક્યાં.
ટાર્સ્કીએ હાર્વર્ડ,
ન્યૂ યૉર્કની સિટી
કૉલેજ, પ્રિન્સ્ટનની
ઇન્સ્ટિટ્યૂટ ફૉર
ઍડવાન્સ્ડ સ્ટડી
અને છેવટે યુનિવર્સિટી
ઑફ કૅલિફોર્નિયા,
બર્કલીમાં અધ્યાપન
કર્યું. ત્યાં તેમણે
તર્કશાસ્ત્ર અને વિજ્ઞાનના
પદ્ધતિશાસ્ત્રનો બહુવિષયક
કાર્યક્રમ સ્થાપ્યો.
26 ઑક્ટોબર 1983ના
રોજ 82 વર્ષની
ઉંમરે તેમનું અવસાન
થયું.

\begin{reading} 
ટાર્સ્કીના જીવન
વિશે વધુ જાણવા
\emph{Alfred Tarski: Life
  and Logic}
જીવનચરિત્ર જુઓ
\citep{Feferman2004}.
તાર્કિક ફલિતતા અને
સત્ય અંગેના તેમના
પાયાના લેખો
\citep{Tarski1983}માં
અંગ્રેજીમાં ઉપલબ્ધ
છે. તેમનાં બધાં
મૂળ લખાણો ચાર
ખંડની શ્રેણીમાં
એકત્ર કરાયાં છે
\citep{Tarski1981}.
\end{reading}

\end{document}
